\chapter{Methodology}
\label{chap:methodology}

This chapter describes the research methods, theoretical frameworks, and evaluation techniques employed to address the gripper configuration
optimization problem. It is organized according to the principal methodological components of the study: problem formulation, candidate evaluation,
optimization, comparative experimentation, and result analysis. Implementation details and experimental procedures are treated separately in the
approach chapter.

\section{Research Design}

The thesis adopts a \emph{computational experimental methodology}. The gripper design problem is approached through mathematical modeling,
simulation-based evaluation, and controlled comparative optimization experiments. This design is appropriate because the objectives governing gripper
quality cannot be expressed in closed analytical form; they depend on geometric interactions between the gripper and the workpiece that must be
assessed computationally. The methodology therefore combines multi-objective and single-objective optimization frameworks, black-box evaluation
procedures, and empirical algorithm comparison to identify high-quality gripper configurations and to characterize algorithm behavior on this class of
problems.

\section{Problem Formulation}
\label{sec:methodology_problem_formulation}

\subsection{Multi-Objective Optimization}
\label{subsec:methodology_moo}

The suction-cup gripper design task is formulated as a multi-objective optimization problem. Let $\mathbf{q} \in \mathcal{Q}$
\nomenclature{$\mathbf{q}$}{Decision vector representing a candidate gripper configuration} \nomenclature{$\mathcal{Q}$}{Admissible search space of
decision vectors} denote the decision vector and let $d$
\nomenclature{$d$}{Number of decision variables} 
denote the number of decision variables. The problem is formulated as
\begin{equation}
    \min_{\mathbf{q} \in \mathcal{Q}}\;
    \mathbf{f}(\mathbf{q}) = \bigl(f_1(\mathbf{q}),\, f_2(\mathbf{q}),\, \dots,\, f_m(\mathbf{q})\bigr),
\end{equation}
where $m$
\nomenclature{$m$}{Number of objective functions}
is the number of objectives, subject to inequality constraints
\begin{equation}
    g_j(\mathbf{q}) \leq 0, \quad j = 1, \dots, n_c,
    \nomenclature{$n_c$}{Number of explicit inequality constraints}
    \nomenclature{$g_j(\mathbf{q})$}{$j$-th inequality constraint function}
\end{equation}
and variable bounds
\begin{equation}
    q_i^{\mathrm{L}} \leq q_i \leq q_i^{\mathrm{U}}, \quad i = 1, \dots, d.
    \nomenclature{$q_i^{\mathrm{L}},\,q_i^{\mathrm{U}}$}{Lower and upper bounds of the $i$-th decision variable}
\end{equation}
Objectives that are naturally maximized are negated to maintain a consistent minimization framework. A multi-objective formulation is chosen because
the governing performance criteria are inherently conflicting: maximizing geometric support area and maximizing grasp diversity impose competing
demands on the suction-cup layout. Rather than collapsing these into a single weighted aggregate, the multi-objective formulation produces a
Pareto-optimal solution set that makes trade-offs explicit and preserves design alternatives.

\subsection{Single-Objective Optimization}
\label{subsec:methodology_soo}

The two-finger gripper design task is formulated as a single-objective optimization problem, since the primary design criterion is a scalar maximal
proximity score. The problem takes the standard form
\begin{equation}
    \min_{\mathbf{q} \in \mathcal{Q}}\; f(\mathbf{q}),
\end{equation}
subject to the same constraint and bound structure defined above.

\subsection{Decision Variables}
\label{subsec:methodology_decision_variables}

For the suction-cup gripper, the decision vector
\begin{equation}
    \mathbf{q} = [r,\, \alpha_1,\, \alpha_2,\, \alpha_3]
\end{equation}
parameterizes the triangular suction-cup layout, where $r$ 
\nomenclature{$r$}{Common radial distance of the suction cups from the gripper center} 
is the common radial distance from the gripper center and $\alpha_i$ \nomenclature{$\alpha_i$}{Angular position of the $i$-th suction cup, $i \in
\{1,2,3\}$} are the angular positions of the three cups. The Cartesian cup positions are recovered via the transformation $\mathbf{p} =
T(\mathbf{q})$, \nomenclature{$T(\mathbf{q})$}{Transformation mapping the reduced decision vector to Cartesian suction-cup positions}
\nomenclature{$\mathbf{p}$}{Cartesian position vector of the suction cups} which maps the reduced parameterization to the configuration used during
evaluation.

For the two-finger gripper, the decision vector
\begin{equation}
    \mathbf{q} = [d_{\mathrm{gn}},\, l_{\mathrm{gn}},\, h_{\mathrm{spline},2},\, h_{\mathrm{spline},3},\, h_{\mathrm{spline},4}]
\end{equation}
controls the finger geometry, where $d_{\mathrm{gn}}$ \nomenclature{$d_{\mathrm{gn}}$}{Gripper nose depth of the two-finger gripper} and
$l_{\mathrm{gn}}$ \nomenclature{$l_{\mathrm{gn}}$}{Gripper nose length of the two-finger gripper} define the nose geometry, and
$h_{\mathrm{spline},i}$ \nomenclature{$h_{\mathrm{spline},i}$}{Vertical position of the $i$-th variable spline control point, $i \in \{2,3,4\}$} for
$i \in \{2,3,4\}$ are the variable control point heights of the finger contour spline.

\subsection{Objective Functions}
\label{subsec:methodology_objectives}

For the suction-cup gripper, two objectives are defined. The first is the support triangle area $A(\mathbf{p})$, \nomenclature{$A(\mathbf{p})$}{Area
of the triangle formed by the three suction-cup positions} which serves as a geometric proxy for grasp stability. The second is a grasp diversity
metric $D(\mathbf{p})$, \nomenclature{$D(\mathbf{p})$}{Grasp diversity metric for suction-cup configuration $\mathbf{p}$} defined as a weighted
combination of spatial and orientational diversity:
\begin{equation}
    D(\mathbf{p}) = w_{\mathrm{s}}\, D_{\mathrm{s}}(\mathbf{p}) + w_{\mathrm{o}}\, D_{\mathrm{o}}(\mathbf{p}), 
    \nomenclature{$D_{\mathrm{s}}(\mathbf{p})$}{Spatial diversity component of the grasp set} 
    \nomenclature{$D_{\mathrm{o}}(\mathbf{p})$}{Orientational diversity component of the grasp set} 
    \nomenclature{$w_{\mathrm{s}},\,w_{\mathrm{o}}$}{Weighting coefficients for the spatial and orientational diversity components} 
\end{equation}
where $D_{\mathrm{s}}$ measures the spatial spread of feasible grasp contact points over the workpiece surface, normalized by a characteristic object
scale, and $D_{\mathrm{o}}$ measures pairwise rotational differences between grasp orientations. Both components are normalized to $[0,1]$. The
objective vector is
\begin{equation}
    \mathbf{f}(\mathbf{q}) = \bigl[-A(\mathbf{p}),\; -D(\mathbf{p})\bigr], \qquad \mathbf{p} = T(\mathbf{q}).
\end{equation}

For the two-finger gripper, the objective is the negated maximal proximity score
\begin{equation}
    f(\mathbf{q}) = -s_{\max}(\mathbf{q}),
    \nomenclature{$s_{\max}(\mathbf{q})$}{Maximal collision-free voxel proximity score
    for two-finger gripper configuration $\mathbf{q}$}
\end{equation}
where $s_{\max}$ is the highest proximity score across all collision-free placements of the gripper relative to the workpiece, as produced by the
voxel-based evaluation described in Section~\ref{subsec:methodology_voxel}. A placement is considered admissible only if the required finger opening
width $w_{\mathrm{open}}(\mathbf{q})$ \nomenclature{$w_{\mathrm{open}}(\mathbf{q})$}{Required finger opening width at the optimal grasp placement;
negative values indicate a valid closing motion} is negative, ensuring that only configurations requiring a physical closing motion are evaluated.
Placements violating this condition are assigned $s_{\max} = 0$.

\section{Evaluation Method}
\label{sec:methodology_evaluation}

\subsection{Evaluation Architecture}
\label{sec:methodology_evaluation_architecture}

The evaluation of candidate gripper configurations is decoupled from the optimization
process through a client--server architecture consisting of three independent pipelines:
an \emph{Optimization Pipeline} and two \emph{Evaluation Pipelines}. This separation
constitutes a methodological design choice that enforces a strict black-box interface
between the optimization strategy and the objective function computation, ensuring that
the optimization algorithms operate solely on input--output pairs
$(\mathbf{q},\, \mathbf{f}(\mathbf{q}))$ without access to internal evaluation logic.
This architecture also makes the framework generalizable: any optimization algorithm
conforming to the interface can be applied to any evaluation pipeline without
modification of either component.

\subsubsection{Optimization Pipeline}
\label{subsec:methodology_opt_pipeline}

The Optimization Pipeline\nomenclature{OP}{Optimization Pipeline; the central
client-side component responsible for executing optimization algorithms and
querying evaluation services} serves as the central orchestration component.
It implements a generalized optimization interface that is agnostic to the specific
gripper type or evaluation method. At each objective function query, the pipeline
acts as a client and issues a request carrying the current decision vector
$\mathbf{q}$ to the appropriate Evaluation Pipeline. The returned objective
value(s) are then passed to the active optimization algorithm. This design
enforces a clean separation between the \emph{search strategy} and the
\emph{objective evaluation}, which is the defining characteristic of black-box
optimization methodology.

\subsubsection{Suction-Cup Evaluation Pipeline}
\label{subsec:methodology_sc_pipeline}

The Suction-Cup Evaluation Pipeline\nomenclature{SCEP}{Suction-Cup Evaluation
Pipeline; the API server responsible for computing suction-cup gripper objectives}
operates as a server that receives decision vectors of the form
$\mathbf{q} = [r,\, \alpha_1,\, \alpha_2,\, \alpha_3]$ from the Optimization
Pipeline client. Upon receipt, it computes the two objective quantities---the
support triangle area $A(\mathbf{p})$ and the grasp diversity metric
$D(\mathbf{p})$---and returns them to the client. The API interface ensures
that the objective computation is encapsulated and that the optimization
algorithm has no visibility into the geometric calculations performed internally.

\subsubsection{Two-Finger Evaluation Pipeline}
\label{subsec:methodology_tf_pipeline}

The Two-Finger Evaluation Pipeline\nomenclature{TFEP}{Two-Finger Evaluation
Pipeline; the API server responsible for computing the maximal proximity score
for two-finger gripper configurations} operates as a server that receives
decision vectors of the form
$\mathbf{q} = [d_{\mathrm{gn}},\, l_{\mathrm{gn}},\,
h_{\mathrm{spline},2},\, h_{\mathrm{spline},3},\, h_{\mathrm{spline},4}]$
from the Optimization Pipeline client. Upon receipt, it computes the maximal
proximity score $s_{\max}(\mathbf{q})$ via the voxel-based evaluation method
described in Section~\ref{subsec:methodology_voxel}, and returns the scalar
value to the client.


\subsection{Voxel-Based Proximity Scoring for the Two-Finger Gripper}
\label{subsec:methodology_voxel}

The two-finger gripper is evaluated using a voxel-based proximity scoring method. Both the gripper and the workpiece are represented as discrete
occupancy grids $G_{\mathrm{occ}}$ \nomenclature{$G_{\mathrm{occ}}$}{Voxelized occupancy grid of the two-finger gripper geometry} and
$O_{\mathrm{occ}}$, \nomenclature{$O_{\mathrm{occ}}$}{Voxelized occupancy grid of the workpiece geometry} with voxel size
$v_s$.\nomenclature{$v_s$}{Voxel size used for discretization of gripper and workpiece geometries} The proximity of the finger surfaces to the
workpiece is modeled by a three-dimensional Gaussian kernel with standard deviations $\sigma_l$, $\sigma_w$, $\sigma_h$ 
\nomenclature{$\sigma_l,\,\sigma_w,\,\sigma_h$}{Standard deviations of the Gaussian proximity kernel in the length, width, and height directions} 
in the length, width, and height directions respectively. The proximity score over all admissible relative placements of the gripper and workpiece is
computed efficiently via the convolution theorem using FFT. Directional cropping of the kernel enforces that each finger contributes proximity only
from its inward-facing side. Placements with gripper-workpiece collisions are assigned a score of $-\infty$ and are excluded. This method is
appropriate because it jointly evaluates all candidate placements in a single computation, avoiding the need for explicit grasp candidate enumeration.

\subsection{Grasp Diversity Metric for the Suction-Cup Gripper}

Grasp diversity is quantified through the metric $D(\mathbf{p})$ defined in Section~\ref{subsec:methodology_objectives}. The spatial component
$D_{\mathrm{s}}$ is based on the pairwise distances between grasp contact points, normalized to a characteristic object scale for cross-workpiece
comparability. The orientational component $D_{\mathrm{o}}$ is based on pairwise geodesic distances between grasp orientations in rotation space. This
two-component formulation captures both positional coverage and orientational variation of the grasp set, which together characterize the practical
flexibility of a gripper configuration.

\section{Optimization Methods}
\label{sec:methodology_optimization}

\subsection{Evolutionary Multi-Objective Algorithms}

For the both grippers, population-based evolutionary algorithms are employed. These methods are well suited to multi-objective black-box
problems because they maintain a population of candidate solutions, approximate the full Pareto front in a single run, and require no gradient
information. The study considers NSGA-II, NSGA-III, MOEA/D, CMOPSO, and SMS-EMOA as representative algorithms covering different Pareto front
approximation strategies, including dominance-based selection, reference-direction decomposition, and hypervolume-indicator guidance. The choice of
algorithm is treated as an empirical question rather than an a priori assumption.

\subsection{Bayesian Optimization}

Bayesian optimization is included as a comparative strategy for both gripper types. It is motivated by its sample efficiency in expensive black-box
settings: a probabilistic surrogate model is maintained over the objective function and an acquisition function is used to guide evaluation toward
promising regions of the search space. For the suction-cup gripper, multi-objective Bayesian optimization is considered; for the two-finger gripper,
single-objective Bayesian optimization is used. The method operates on the same problem definition as the evolutionary algorithms, ensuring that
performance differences are attributable solely to the optimization strategy.

\subsection{Termination and Hyperparameter Treatment}

Optimization runs are terminated based on a combination of a maximum number of function evaluations, a maximum number of generations, and a
convergence criterion evaluated over a sliding window in objective space. This compound stopping rule prevents premature termination while avoiding
unnecessary computation after convergence. Algorithm hyperparameters are selected through a structured tuning procedure based on repeated runs
evaluated by the hypervolume indicator, reducing the risk of performance bias from incidental parameter choices.

\section{Comparative Experimental Method}
\label{sec:methodology_comparative_experiments}

A controlled comparison protocol is adopted to ensure that observed performance differences are attributable to the optimization methods rather than
to differences in experimental conditions. All algorithms are evaluated on the same workpiece data, gripper parameterization, objective functions,
feasibility conditions, evaluation pipeline, and evaluation budget. Each algorithm is executed under multiple independent random seeds to assess not
only mean performance but also variability and robustness across runs.

\section{Performance Indicators and Result Analysis}
\label{sec:methodology_analysis}

Algorithm performance is assessed using standard multi-objective and single-objective indicators. For the suction-cup gripper, the primary indicator
is the \emph{hypervolume indicator}, which measures the volume of objective space dominated by the solution set with respect to a reference point,
thereby jointly capturing convergence quality and front diversity. For the two-finger gripper, the primary performance indicator is the best achieved
maximal proximity score $s_{\max}$ and its convergence profile over function evaluations. Robustness is characterized by the distribution of
$s_{\max}$ across repeated independent runs.